\clearpage
\hypertarget{native-475-lower}{}
\begin{center}
\rule{0.55\linewidth}{0.4pt}\\[2pt]
Astronomische Nachrichten Vol. XXII. No. 526. pp. 353. 354. 25 January 1845.\\[2pt]
\rule{0.55\linewidth}{0.4pt}
\end{center}

On the 11th, Mr \textsc{Petersen} compared the comet with two stars from the Hist. Cél., pp. 237 and 295, for which he assumed the following apparent positions at the time of observation

\begin{center}
\begin{tabular}{rr}
$19^{\rm u}\ 21^{\rm m}\ 2^{\rm s}85$ & $+42^\circ\ 1'\ 5''1$ \\
$22\ 20.95$ & $41\ 55\ 11.8$ \\
\end{tabular}
\end{center}
\noindent He found:

\begin{center}
Jan. 11.\quad $9^{\rm u}\ 16^{\rm m}\ 7^{\rm s}$ Altona mean time\quad $289^\circ\ 37'\ 35''6$\quad $+41^\circ\ 58'\ 40''7$.
\end{center}

From two comparisons with the same stars, Councillor \textsc{Gauss} found

\begin{center}
\begin{tabular}{lll}
Jan. 11. & $5^{\rm u}\ 35^{\rm m}\ 19^{\rm s}2$ Göttingen mean time & $+41^\circ\ 54'\ 40''2$ \\
 & $5\ 39\ 29.0$ & R.A. $289^\circ\ 41'\ 47''6$ \\
\end{tabular}
\end{center}

The declination is based on the second star alone and the right ascension on the mean of the two. Councillor \textsc{Gauss} would like the positions of the stars now to be determined at lower culmination.

As I finish this article, I receive another Göttingen observation from Councillor \textsc{Gauss}, the following elements from Mr \textsc{Rümker}, and an improvement of his earlier orbit from Mr \textsc{Sievers}. Mr \textsc{Rümker} would have supplied his elements sooner had his duties at the Navigation School, where he could use more help, not taken up almost all his time at present.

\begin{center}
\setlength{\tabcolsep}{3pt}
\begin{tabular}{lll}
1845 Jan. 12. & $5^{\rm u}\ 51^{\rm m}\ 11^{\rm s}0$ Göttingen mean time & R.A. $289^\circ\ 14'\ 45''5$ \\
 & $5\ 54\ 59.5$ & Decl. $42\ 22\ 15.4$ N. \\
\end{tabular}
\end{center}

\Needspace{7\baselineskip}
\SourceLeaf{476}
For comparison, 3 stars were used. Their positions for the day of observation, derived from the Hist. Cél., are as follows:

\begin{center}
\begin{tabular}{rr}
$288^\circ\ 54'\ 5''4$ & $+42^\circ\ 15'\ 49''7$ \\
$289\ 1\ 5.6$ & $42\ 23\ 59.7$ \\
$290\ 7\ 35.2$ & $42\ 29\ 48.2$ \\
\end{tabular}
\end{center}

The last of these stars is plotted considerably too far north on \textsc{Harding}’s chart.
\begin{flushright}\textsc{Schumacher}.\end{flushright}
